\documentclass[12pt,a4paper]{article}
%
%	STARWARE - Stile base per la scrittura di documenti interni in LaTex
%

%
%	Pacchetti globali
%	Fare qui eventuali aggiunte!
%
\usepackage[utf8]{inputenc}
\usepackage[italian]{babel}
\usepackage[babel]{csquotes}
\usepackage{url}
\usepackage{graphicx}
\usepackage[colorlinks]{hyperref}
\usepackage{lastpage}
\usepackage{fancyhdr}
\usepackage[top=1cm,bottom=4cm,left=80pt,right=80pt]{geometry} %disegna la linea
\usepackage{listings} %per grandi porzioni di codice
\usepackage{color}
\usepackage[table]{xcolor}
\usepackage{booktabs,tabularx}
\usepackage{makeidx}
\usepackage{fixltx2e}
\usepackage{hyperref}
\usepackage{enumitem}
\usepackage{color}
\usepackage[T1]{fontenc}
\usepackage{float}
\usepackage{svg}
\usepackage{amsmath}
\usepackage[toc]{glossaries}
\usepackage{dirtree}
\usepackage{listings}

\makeglossaries

\bibliographystyle{alpha}

%
%	VARIABILI GLOBALI
%
\newcommand{\nomeGruppo}{StarWare}
\newcommand{\mailGruppo}{starware.swe@gmail.com}
\newcommand{\uni}{Universit\`{a} degli Studi di Padova}
\newcommand{\uniAA}{2015/2016}
\newcommand{\Cardin}{Prof. Riccardo Cardin}
\newcommand{\Vardanega}{Prof. Tullio Vardanega}
\newcommand{\Zucchetti}{Zucchetti S.p.a.}
\newcommand{\prj}{Quizzipedia}
\newcommand{\prjL}{Quizzipedia: software per la gestione di questionari}

\newcommand{\AVI}{Alessio Vitella}
\newcommand{\AVE}{Andrea Venier}
\newcommand{\NDC}{Nicola De Cao}
\newcommand{\IB}{Igor Baylyak}
\newcommand{\WS}{Walter Sandon}
\newcommand{\TP}{Thomas Pigarelli}
\newcommand{\AB}{Anna Bonaldo}

\newcommand{\mgls}[1]{\gls{#1}\textsubscript{G}}
\newcommand{\mglspl}[1]{\glspl{#1}\textsubscript{G}}
\newcommand{\mGls}[1]{\Gls{#1}\textsubscript{G}}
\newcommand{\mGlspl}[1]{\Glspl{#1}\textsubscript{G}}

\newcommand{\AM}{\emph{\mGls{amministratore}}}
\newcommand{\AN}{\emph{\mGls{analista}}}
\newcommand{\PG}{\emph{\mGls{progettista}}}
\newcommand{\PR}{\emph{\mGls{programmatore}}}
\newcommand{\VR}{\emph{\mGls{verificatore}}}
\newcommand{\PM}{\emph{\mGls{project manager}}}

\newcommand{\AMpl}{\emph{\mGlspl{amministratore}}}
\newcommand{\ANpl}{\emph{\mGlspl{analista}}}
\newcommand{\PGpl}{\emph{\mGlspl{progettista}}}
\newcommand{\PRpl}{\emph{\mGlspl{programmatore}}}
\newcommand{\VRpl}{\emph{\mGlspl{verificatore}}}
\newcommand{\PMpl}{\emph{\mGlspl{project manager}}}

\newcommand{\RR}{\emph{\mGls{revisione dei requisiti}}}
\newcommand{\RA}{\emph{\mGls{revisione di accettazione}}}
\newcommand{\RP}{\emph{\mGls{revisione di progettazione}}}
\newcommand{\RQ}{\emph{\mGls{revisione di qualifica}}}

\newcommand{\NdP}{\emph{\mGls{norme di progetto}}}
\newcommand{\SdF}{\emph{\mGls{stidio di fattibilita}}}
\newcommand{\AdR}{\emph{\mGls{analisi dei requisiti}}}
\newcommand{\PdP}{\emph{\mGls{piano di progetto}}}
\newcommand{\PdQ}{\emph{\mGls{piano di qualifica}}}

\newcommand{\latex}[1]{\texttt{#1}}
\newcommand{\fileName}[1]{\texttt{#1}}
\newcommand{\filePath}[1]{\texttt{#1}}
\newcommand{\TODO}[1]{\texttt{\large \color{red} \underline{TODO: #1}}}

\newcommand{\licenza}{GNU GENERAL PUBLIC LICENSE V2}

%per compilare il template usare questi sotto e commentare glia altri:
%\newcommand{\logoLungo}{../imgs/logoLungo.png}
%\newcommand{\logoGrande}{../imgs/logoGrande.png}
\newcommand{\logoLungo}{../../../template/imgs/logoLungo.png}
\newcommand{\logoGrande}{../../../template/imgs/logoGrande.png}

%
%	Setup stili
%

\newcommand{\HRule}{\rule{\linewidth}{0.5mm}}

\definecolor{dkgreen}{rgb}{0,0.6,0}
\definecolor{gray}{rgb}{0.5,0.5,0.5}
\definecolor{mauve}{rgb}{0.58,0,0.82}
\definecolor{light}{RGB}{255,255,190}

%
%	Setup di pagina
%

%colorazione link
\hypersetup
{
	colorlinks=true,
	linkcolor=black,
	urlcolor=blue,
	citecolor=blue
}

%	Setup Header + Footer

\pagestyle{fancy}
\setlength{\headheight}{2cm} %settato grandezza header

\renewcommand{\footrulewidth}{0.5pt} %ridefinisco il valore della riga di intestazione
\renewcommand{\headrulewidth}{0.5pt} %ridefinisco il valore della riga di pie' di pagina
\addtolength{\headwidth}{\marginparsep}
\addtolength{\headwidth}{\marginparwidth}

\fancyhead{} %annulla head di default
\fancyfoot{} %annulla foot di default

%	Logo intestazione
\lhead{\includegraphics[scale=0.12]{\logoLungo}}


%	footer
\cfoot{
	\uni \ - \uniAA \\
	\href{mailto:\mailGruppo}{\mailGruppo}\\
	{\tiny Questo documento è distribuito sotto licenza {\licenza}}
}
\rfoot{
	\thepage\ di \pageref{LastPage}
}


% test subsubsubsection

\usepackage{titlesec}
\usepackage{hyperref}

\titleclass{\subsubsubsection}{straight}[\subsection]

\newcounter{subsubsubsection}[subsubsection]
\renewcommand\thesubsubsubsection{\thesubsubsection.\arabic{subsubsubsection}}
\renewcommand\theparagraph{\thesubsubsubsection.\arabic{paragraph}} % optional; useful if paragraphs are to be numbered

\titleformat{\subsubsubsection}
  {\normalfont\normalsize\bfseries}{\thesubsubsubsection}{1em}{}
\titlespacing*{\subsubsubsection}
{0pt}{3.25ex plus 1ex minus .2ex}{1.5ex plus .2ex}

\makeatletter
\renewcommand\paragraph{\@startsection{paragraph}{5}{\z@}%
  {3.25ex \@plus1ex \@minus.2ex}%
  {-1em}%
  {\normalfont\normalsize\bfseries}}
\renewcommand\subparagraph{\@startsection{subparagraph}{6}{\parindent}%
  {3.25ex \@plus1ex \@minus .2ex}%
  {-1em}%
  {\normalfont\normalsize\bfseries}}
\def\toclevel@subsubsubsection{4}
\def\toclevel@paragraph{5}
\def\toclevel@paragraph{6}
\def\l@subsubsubsection{\@dottedtocline{4}{7em}{4em}}
\def\l@paragraph{\@dottedtocline{5}{10em}{5em}}
\def\l@subparagraph{\@dottedtocline{6}{14em}{6em}}
\makeatother

\setcounter{secnumdepth}{4}
\setcounter{tocdepth}{4}


%pdflatex -synctex=1 -interaction=nonstopmode %.tex|makeglossaries %|pdflatex -synctex=1 -interaction=nonstopmode %.tex|pdflatex -synctex=1 -interaction=nonstopmode %.tex
\makeglossary

\newglossaryentry{responsabile} {
	name=responsabile,
	description={è il responsabile della gestione, pianificazione e realizzazione del progetto},
	plural=Responsabili
}

\newglossaryentry{verificatore} {
	name=verificatore,
	description={è il responsabile dell'attività di verifica},
	plural=Verificatori
}

\newglossaryentry{programmatore} {
	name=programmatore,
	description={è responsabile delle attività di codifica miranti alla realizzazione del prodotto e delle componenti di ausilio necessarie per l'esecuzione delle prove di verifica e validazione},
	plural=programmatori
}

\newglossaryentry{progettista} {
	name=progettista,
	description={è responsabile delle attività di progettazione},
	plural=Progettisti
}

\newglossaryentry{analista} {
	name=analista,
	description={è responsabile delle attività di analisi. },
	plural=Analisti
}

\newglossaryentry{amministratore} {
	name=amministratore,
	description={è responsabile dell'efficienza e dell'operatività dell'ambiente di sviluppo; si occupa della redazione e attuazione di piani e procedure di gestione della qualità; inoltre gestisce l'archivio della documentazione del progetto},
	plural=Amministratori
}

\newglossaryentry{revisione dei requisiti} {
	name=revisione dei requisiti,
	description={è una revisione formale che determina l'accesso del gruppo al progetto didattico e la concordanza con il cliente di una visione condivisa del prodotto atteso}
}

\newglossaryentry{revisione di accettazione} {
	name=revisione di accettazione,
	description={è una revisione formale per l'accertamento del soddisfacimento di tutti i requisiti e il completamento del progetto}
}

\newglossaryentry{revisione di progettazione} {
	name=revisione di progettazione,
	description={è una revisione di progresso che accerta la realizzabilità del prodotto e informa il cliente sulle caratteristiche del prodotto}
}

\newglossaryentry{revisione di qualifica} {
	name=revisione di qualifica,
	description={è una revisione di progresso che approva l'esito finale delle verifiche e attiva la fase di validazione}
}

\newglossaryentry{analisi} {
	name=analisi,
	description={è il periodo di preparazione e produzioni di documenti che precede la Revisione dei requisiti}
}

\newglossaryentry{progettazione} {
	name=progettazione,
	description={è il periodo che intercorre tra la Revisione dei requisiti e la Revisione di progettazione}
}

\newglossaryentry{codifica} {
	name=codifica,
	description={è il periodo che intercorre tra la Revisione di progettazione e la Revisione di qualifica}
}

\newglossaryentry{validazione} {
	name=validazione,
	description={è il periodo che intercorre tra la Revisione di qualifica e la Revisione di accettazione}
}

\newglossaryentry{repository} {
	name=repository,
	description={è dove i file sono memorizzati, spesso su un server}
}

\newglossaryentry{ruolo} {
	name=ruolo,
	description={una delle figure professionali che una persona fisica interpreta nel corso del progetto. I ruoli sono: responsabile, amministratore, analista, progettista, programmatore e verificatore},
    plural=ruoli
}

\newglossaryentry{svg} {
	name=SVG,
	description={è un formato per la visualizzazione di oggetti in grafica vettoriale. Per maggiori informazioni si veda \href{https://it.wikipedia.org/wiki/Scalable_Vector_Graphics}{qui}}
}

\newglossaryentry{png} {
	name=PNG,
	description={abbreviazione di Portable Network Graphics, è un formato di file per memorizzare immagini. Per ulteriori informazioni si veda \href{http://it.wikipedia.org/wiki/Portable_Network_Graphics}{qui}}
}

\newglossaryentry{pdf} {
	name=PDF,
	description={è un formato di file basato su un linguaggio di descrizione di pagina sviluppato da Adobe Systems nel 1993 per rappresentare documenti in modo indipendente dall’hardware e dal software utilizzati per generarli o per visualizzarli. Per ulteriori informazioni si veda \href{http://it.wikipedia.org/wiki/Portable_Document_Format}{qui}}
}

\newglossaryentry{uml} {
	name=UML,
	description={è un linguaggio di modellazione e specifica basato sul paradigma object-oriented. Per ulteriori informazioni si veda \href{http://it.wikipedia.org/wiki/Unified_Modeling_Language}{qui}}
}

\newglossaryentry{walkthrough} {
	name=walkthrough,
    description={consiste nella lettura di un documento o codice cercando errori ed anomalie senza un'idea precisa di quali tipi di errori sarà possibile trovare}
}

\newglossaryentry{lista di controllo} {
	name=lista di controllo,
	description={è un elenco di cose da fare per eseguire una determinata attività}
}

\newglossaryentry{inspection} {
	name=inspection,
	description={è la lettura mirata di un documento o codice cercando errori specifici}
}

\newglossaryentry{milestone} {
	name=milestone,
	description={momento saliente nello sviluppo di un prodotto software per la quale devono essere pronti documenti e/o funzionalità}
}

\newglossaryentry{ticket} {
	name=ticket,
	description={rappresenta un compito nell'organizzazione e distribuzione del lavoro all'interno del progetto},
	plural=tickets
}

\newglossaryentry{commit} {
	name=commit,
	description={è la copia di modifiche fatte su file locali verso la repository remota. Esso rappresenta anche un particolare stato della repository nel tempo}
}

\newglossaryentry{versionamento} {
	name=versionamento,
	description={è la gestione di un versioni multiple di un insieme di informazioni. Per maggiori informazioni si veda \href{http://it. wikipedia.org/wiki/Controllo_versione}{qui}}
}

\newglossaryentry{task} {
	name=task,
	description={è un compito secondo la definizione dello standard IEEE 12207},
	plural=tasks
}

\newglossaryentry{attivita} {
	name=attivita,
	description={è un insieme di task}
}

\newglossaryentry{redattore} {
	name=redattore,
	description={colui che redige un documento},
	plural=redattori
}

\newglossaryentry{proponente} {
	name=proponente,
	description={colui che ha proposto al committente un capitolato d'appalto}
}

\newglossaryentry{committente} {
	name=committente,
	description={colui che assegna un compito. In questo caso è il Professor Tullio Vardanega}
}

\newglossaryentry{quality assurance} {
	name=quality assurance,
	description={è l'insieme delle attività volte a garantire il soddisfacimento degli obiettivi della qualità}
}

\newglossaryentry{telegram} {
	name=telegram,
	description={è un servizio di messaggistica istantanea utilizzato dal gruppo per comunicazioni interne. Per maggiori informazioni si veda \href{https://it.wikipedia.org/wiki/Telegram_(software)}{qui}}
}

\newglossaryentry{browser} {
	name=browser,
	description={è un'applicazione per il recupero, la presentazione e la navigazione di risorse web}
}

\newglossaryentry{google drive} {
	name=Google Drive,
	description={è un servizio di memorizzazione e sincronizzazione online introdotto da Google il 24 aprile 2012. Per maggiori informazioni si veda \href{https://it.wikipedia.org/wiki/Google_Drive}{qui}}
}

\newglossaryentry{skype} {
	name=Skype,
	description={è un software proprietario freeware di messaggistica istantanea e VoIP. Per maggiori informazioni si veda \href{https://it.wikipedia.org/wiki/Skype}{qui}}
}

\newglossaryentry{gantt} {
	name=Gantt,
	description={è un diagramma di supporto alla gestione dei progetti}
}

\newglossaryentry{projectlibre} {
	name=ProjectLibre,
	description={è un software di gestione progettuale}
}

\newglossaryentry{pert} {
	name=PERT,
	description={è uno strumento volto alla programmazione delle attività che compongono il progetto e, più in generale, alla gestione degli aspetti temporali di quest'ultimo}
}

\newglossaryentry{subtask} {
	name=subtask,
	description={è un task compreso all'interno di un altro task. La totalità di tutti i subtasks costituisce un intero task}
	plural=subtasks
}

\newglossaryentry{ticketing} {
	name=ticketing,
	description={procedura con la quale il Responsabile assegna un task}
}

\newglossaryentry{git} {
	name=git,
	description={è un sistema software di controllo di versione distribuito}
}

\newglossaryentry{quizzpedia} {
	name=Quizzpedia,
	description={è il nome del prodotto software richiesto dal capitolato d'appalto scelto}
}

\newglossaryentry{schierabile} {
	name=schierabile,
	description={è la capacità di rilasciare al cliente, con relativa installazione e messa in funzione o esercizio, di una applicazione o di un sistema software tipicamente all'interno di un sistema informatico aziendale}
}

\newglossaryentry{cross-platform} {
	name=cross-platform,
	description={può essere riferito ad un linguaggio di programmazione, ad un'applicazione software o ad un dispositivo hardware che funziona su più di un sistema}
}

\newglossaryentry{qml} {
	name=QML,
	description={è un "Domain Specific Language" richiesto dal capitolato d'appalto per la definizione delle domande all'interno del sistema}
}

\newglossaryentry{tomcat} {
	name=Tomcat,
	description={è un application server nella forma di contenitore servlet open source sviluppato dalla Apache Software Foundation. Per maggiori informazioni si veda \href{https://it.wikipedia.org/wiki/Apache_Tomcat}{qui}}
}

\newglossaryentry{java} {
	name=Java,
	description={è un linguaggio di programmazione orientato agli oggetti, specificatamente progettato per essere il più possibile indipendente dalla piattaforma di esecuzione. Per maggiori informazioni si veda \href{https://it.wikipedia.org/wiki/Java_(linguaggio_di_programmazione)}{qui}}
}

\newglossaryentry{node.js} {
	name=Node.js,
	description={è un framework che permette di realizzare web application usando un linguaggio di programmazione che utilizza la stessa sintassi di JavaScrip. Utilizza un modello event-driven, anzichè il classico modello a processi o thread concorrenti, e ciò significa che si eseguono azioni solo al verificarsi di un evento. Questo modello asincrono rende leggero ed efficiente, ideale per applicazioni real-time in per dispositivi distribuiti}
}

\newglossaryentry{javascript} {
	name=JavaScript,
	description={linguaggio di scripting orientato agli oggetti comunemente usato nella programmazione
Web}
}

\newglossaryentry{postgresql} {
	name=PostgreSQL,
	description={e un sistema di gestione di basi di dati open source usato per applicazioni che richiedono caratteristiche molto complesse complesse}
}

\newglossaryentry{mongodb} {
	name=MongoDB,
	description={è un sistema gestionale di basi di dati NoSQL orientato ai documenti, adatto per ambienti che hanno la necessità d'immagazzinare grosse quantità di dati, e dove il linguaggio utilizzato per la gestione dei dati è \mgls{javascript}}
}

\newglossaryentry{html5} {
	name=HTML5,
	description={linguaggio di markup per la strutturazione delle pagine web}
}

\newglossaryentry{css3} {
	name=CSS3,
	description={è un linguaggio di programmazione web utilizzato per descrivere l'aspetto e la formattazione di un sito web al browser lato client. }
}

\newglossaryentry{xml} {
	name=XML,
	description={è un meta-linguaggio che fornisce un insieme standard di regole sintattiche per modellare la struttura di documenti e dati. Questo insieme di regole, definiscono le modalità secondo cui è possibile crearsi un proprio linguaggio di markup}
}

\newglossaryentry{nosql} {
	name=NoSQL,
	description={come dice anche il termine questi database non sono basati su SQL, non sono basati su uno schema relazionale. I database relazionali sono infatti ottimi quando esistono delle relazioni tra i dati che salviamo, ma sono poco performanti nel caso sia necessario salvare una grande quantità di dati, magari usando la scalabilità orizzontale, quando cioè si utilizzano più server dove salvare questi dati e non solamente incrementando la potenza di un singolo server.
}
}

\newglossaryentry{scala} {
	name=Scala,
	description={è un linguaggio di programmazione di tipo general-purpose multi-paradigma studiato per integrare le caratteristiche e funzionalità dei linguaggi orientati agli oggetti e dei linguaggi funzionali. Per maggiori informazioni si veda \href{https://it.wikipedia.org/wiki/Scala_(linguaggio_di_programmazione)}{qui}}
}

\newglossaryentry{akka} {
	name=Akka,
	description={è un toolkit di strumenti per la costruzione di applicazioni con elevata concorrenza di dati che necessitano di un sistema resilente per l'invio e la ricezione di messaggi}
}

\newglossaryentry{ble} {
	name=BLE,
	description={Bluetooth low energy,  pur mantenendo un range di comunicazione simile a quello classico, fornisce un  consumo energetico dei device notevolmente ridotto}
}

\newglossaryentry{mqtt} {
	name=MQTT,
	description={è un protocollo di messaggistica leggero posizionato in cima a TCP/IP, disegnato per le situazioni in cui è richiesto un basso impatto e dove la banda è limitata. }
}

\newglossaryentry{aws} {
	name=AWS,
	description={è un insieme di servizi di elaborazione rermoti, detti anche servizi web, che costituiscono una piattaforma di cloud computing offerto da Amazon. Per maggiori informazioni si veda \href{https://aws.amazon.com/it/}{qui}}
}

\newglossaryentry{heroku} {
	name=Heroku,
	description={è una  delle prime cloud Platform-as-a-Service (PaaS) che supportava solo Java come linguaggio di programmazione, oggi giorno ha aggiunto molti altri linguaggi come Scala, Phyton, etc. Per maggiori informazioni si veda \href{https://www.heroku.com}{qui}}
}

\newglossaryentry{github} {
	name=Github,
	description={software di controllo di versione che permette di aggiornare un file senza dover sovrascrivere le versioni precedenti}
}

\newglossaryentry{template} {
	name=template,
	description={traducibile in italiano come modello, indica o un programma o un documento idealizzato come un documento semicompilato cartaceo che ha degli spazi bianchi che saranno successivamente riempiti}
	plural=templates
}

\newglossaryentry{dropbox} {
	name=Dropbox,
	description={è un software di cloud storage multipiattaforma, che offre un servizio di file hosting e sincronizzazione automatica di file tramite web}
}

\newglossaryentry{open-source} {
	name=open-source,
	description={un software di cui gli autori ovvero i detentori dei diritti rendono pubblico il codice sorgente, favorendone il libero studio e permettendo a programmatori indipendenti di apportarvi modifiche ed estensioni. Questo è realizzato tramite apposite licenze d'uso}
}

\newglossaryentry{project management} {
	name=project management,
	description={si intende l'insieme di attività aziendali, svolte tipicamene da una figura dedicata e specializzata detta project manager, volte all'analisi, progettazione, pianificazione e realizzazione degli obiettivi di un progetto, gestendolo in tutte le sue caratteristiche e fasi evolutive, nel rispetto di precisi vincoli come i tempi, costi, risorse, scopi, qualità}
}

\newglossaryentry{linux} {
	name=Linux,
	description={è una famiglia di sistemi operativi open-source di tipo Unix-like, rilasciati sotto varie possibili distribuzioni, aventi la caratteristica comune di utilizzare come nucleo il kernel Linux}
}

\newglossaryentry{windows} {
	name=Windows,
	description={è una famiglia di ambienti operativi e sistemi operativi dedicati ai personal computer, alle workstation, ai server e agli smartphone. Il sistema operativo si chiama così per via della sua interfaccia di programmazione di un'applicazione a finestre}
}

\newglossaryentry{mac os} {
	name=Mac OS,
	description={il sistema operativo di Apple dedicato dedicati ai personal computer Macintosh, alle workstation, ai server e agli smartphone  }
}


\newglossaryentry{schedule variance} {
	name=schedule variance,
	description={ogni deviazione alle baseline di un progetto, misurata confrontando costo preventivato di programma di lavoro con il costo preventivato del lavoro svolto. Indica quindi al cliente se il progetto sta procedendo nei tempi stabiliti}
}

\newglossaryentry{cost variance} {
	name=cost variance,
	description={indica al management aziendale se il valore del costo realmente maturato è maggiore, uguale o minore rispetto al costo pianificato}
}


\newglossaryentry{merge} {
	name=merge,
	description={unire due o più quantità}
}

\newglossaryentry{slack} {
	name=slack,
	description={intervallo di tempo entro cui un evento deve avvenire nel rispetto dei vincoli logici e imposti dal reticolo di pianificazione senza compromettere la durata complessiva del progetto}
}

\newglossaryentry{baseline} {
	name=baseline,
	description={di progetto costituisce il punto di riferimento rispetto al quale calcolare gli scostamenti delle principali variabili implicate nella gestione di un progetto}
}

\newglossaryentry{asana} {
	name=Asana,
	description={è un software che permette a dei team di monitorare il loro lavoro e tener traccia dei risultati tramite l'assegnazione di task a componemti specifiche del gruppo di lavoro}
}

\newglossaryentry{deadline} {
	name=deadline,
	description={La deadline di un \mgls{task} è un indicazione dell'urgenza del \mgls{task}; rappresenta un un punto su una linea temporale ideale. Data di scadenza o termine entro il quale deve essere completato un compito assegnato.}
}

\newglossaryentry{revert} {
	name=revert,
	description={per annullare le ultime modifiche effettuate al repository remoto. Per maggiori informazioni si veda \href{https://git-scm.com/docs/}{qui}}
}

\newglossaryentry{backup} {
	name=backup,
	description={si indica la replicazione, su un qualunque supporto di memorizzazione, di materiale informativo archiviato nella memoria di massa, al fine di prevenire la perdita definitiva dei dati in caso di eventi malevoli accidentali o intenzionali}
}

\newglossaryentry{push} {
	name=push,
	description={per inviare modifiche di un documento site in un host locale al repository remoto. Per maggiori informazioni si veda \href{https://git-scm.com/docs/}{qui}}
}

\newglossaryentry{teamwork} {
	name=teamwork,
	description={è un software che permette a dei team di monitorare il loro lavoro e tener traccia dei risultati tramite l'assegnazione di task a componenti specifiche del gruppo di lavoro}
}

\newglossaryentry{evento} {
	name=evento,
	description={\mgls{teamwork}, attraverso il suo calendario, offre la possibilità di impostare eventi in giorni stabiliti dagli utenti. Questi eventi verranno periodicamente segnalati come promemoria a tutte le persone invitate agli stessi.}
}

\newglossaryentry{etichetta} {
	name=etichetta,
	description={è un controllo grafico che mostra informazioni testuali all'interno di un form}
	plural=etichette
}

\newglossaryentry{bug} {
	name=bug,
	description={identifica un errore nella scrittura di un programma software che ne causa un comportamento imprevisto o comunque diverso da quello specificato dal produttore}
	plural=bugs
}

\newglossaryentry{software} {
	name=software,
	description={e’ un termine generico che definisce programmi e procedure utilizzati per far eseguire al computer un determinato compito}
}

\newglossaryentry{desktop} {
	name=desktop,
	description={area dello schermo su cui appaiono le icone e le finestre rappresentanti le memorie di massa collegate al computer ed il loro contenuto}
}

\newglossaryentry{draw.io} {
	name=draw.io,
	description={é un software gratuito per la creazione di diagrammi di flusso, di processo, UML, e diagrammi di rete}
}

\newglossaryentry{tracy} {
	name=tracy,
	description={software opensource per il tracciamento}
}

\newglossaryentry{package} {
	name=package,
	description={è un meccanismo per organizzare classi Java in gruppi logici, principalmente allo scopo di definire namespace distinti per diversi contesti}
}

\newglossaryentry{stakeholder} {
	name=stakeholder,
	description={si indica genericamente un soggetto o un gruppo di soggetti influenti nei confronti di un'iniziativa economica, sia essa un'azienda o un progetto}
}

\newglossaryentry{branch} {
	name=branch,
	description={quando si vuole creare un nuovo ramo al repository remoto. Per maggiori informazioni si veda \href{https://git-scm.com/docs/}{qui}}
	plural=branches
}

\newglossaryentry{pull} {
	name=pull,
	description={un comando gitub per poter ricevere nell'host locale tutte le modifiche fatte nel repository remoto. Per maggiori informazioni si veda \href{https://git-scm.com/docs/}{qui}}
}

\newglossaryentry{underscore} {
	name=underscore,
	description={è un carattere che identifica il trattino basso}
}

\newglossaryentry{spelling} {
	name=spelling,
	description={è l'atto di pronunciare le parole lentamente, separando le singole lettere o le sillabe}
}

\newglossaryentry{cloud} {
	name=cloud,
	description={si indica un sistema di erogazione di risorse informatiche, come l'archiviazione, l'elaborazione o la trasmissione di dati, caratterizzato dalla disponibilità on demand attraverso Internet}
}

\newglossaryentry{smartphone} {
	name=smartphone,
	description={è un telefono cellulare con capacità di calcolo, di memoria e di connessione dati molto più avanzate rispetto ai normali telefoni cellulari}
	plural=smartphones
}

\newglossaryentry{checkbox} {
	name=checkbox,
	description={è un controllo grafico con cui l'utente può effettuare selezioni multiple}
}


\newglossaryentry{makefile} {
	name=makefile,
	description={è usata soprattutto per la compilazione di codice sorgente in codice oggetto, unendo e poi linkando il codice oggetto in programmi eseguibili o in librerie}
}

\newglossaryentry{gulpease} {
	name=Gulpease,
	description={è un indice di leggibilità di un testo tarato sulla lingua italiana. Rispetto ad altri ha il vantaggio di utilizzare la lunghezza delle parole in lettere anziché in sillabe, semplificandone il calcolo automatico}
}

\newglossaryentry{pdca} {
	name=PDCA,
	description={\TODO{}}
}

\usepackage{tikz}
%Titolo documento
\newcommand{\titoloDocumento}{Analisi dei requisiti}

%Prima data di creazione del documento
\newcommand{\dataCreazione}{30 Novembre 2015}

%Inserite la versione attuale del documento
\newcommand{\versione}{1.0.0}

%Stato in cui si trova il documento: Formale solo all'atto di consegna
\newcommand{\stato}{Informale}

%Uso del documento
\newcommand{\uso}{Esterno}

\rhead{\titoloDocumento}
\lfoot{Versione: \versione}
\title{\titoloDocumento}

\begin{document}
\begin{titlepage}
\begin{center}
\today \\
\vspace{1cm}
\begin{Huge}
\textbf{\nomeGruppo} \\
\end{Huge}
\textbf{\prjL} \\
\vspace{1cm}
\includegraphics[scale=0.5]{\logoGrande}
\vspace{1cm}

\HRule \\[0.4cm]
\begin{Huge}
{\huge \bfseries \titoloDocumento}\\[0.4cm]
\end{Huge}
\HRule \\[1cm]
\vfill

\begin{table}[h]
\begin{center}
\begin{tabular}{r | l}
\multicolumn{2}{c}{\textbf{Informazioni sul documento}}\\
\midrule
\textbf{Nome Documento} & \titoloDocumento \\
\textbf{Versione} & \versione \\
\textbf{Stato} & \emph{\stato} \\
\textbf{Uso} & \emph{\uso} \\
\textbf{Data Creazione} & \dataCreazione \\
\textbf{Data Ultima Modifica} & \today \\
\textbf{Redazione} & \IB\\
\  & \TP\\
\  &  \AVE\\
\textbf{Verifica} &  \NDC\\
\ & \AVI \\
\textbf{Approvazione} &  \WS\\
\textbf{Lista Distribuzione} & \nomeGruppo \\
\  & \Vardanega \\
\  & \Cardin \\
\  & Il proponente \Zucchetti \\

\end{tabular}
\end{center}
\end{table}

\end{center}
\end{titlepage}
\newpage

\Large{\textbf{Registro delle modifiche}}\\
\normalsize

\begin{table}[h]
\begin{center}

\begin{tabular}{p{0.12\textwidth} p{0.25\textwidth} p{0.18\textwidth} p{0.4\textwidth}}
\toprule
\textbf{Versione} & \textbf{Autore} & \textbf{Data} & \textbf{Descrizione}\\
\midrule
\midrule

\bottomrule
\end{tabular}
\caption{Versionamento del documento}
\label{tabVers1}
\end{center}
\end{table}
\newpage

\tableofcontents
\listoftables
\newpage

\section{Requisiti}
\begin{longtable}{|r l|p{2cm}|p{6cm}|p{2cm}|}
	\hline
	\multicolumn{2}{|c|}{Requisito} & Tipologia & Descrizione & Fonti\tabularnewline
	\hline
	& \hypertarget{R-3V1}{R-3V1} & Vincolo
	
	Obbligatorio & La parte destinata ai creatori di domande e di questionari dovrà essere fruibile attraverso un PC & Capitolato\tabularnewline
	\hline
	& \hypertarget{R-3V2}{R-3V2} & Vincolo
	
	Obbligatorio & La parte server deve essere realizzata in Java con server Tomcat o Javascaript con server Node.js
	& Capitolato\tabularnewline
	\hline
	& \hypertarget{R-3V3}{R-3V3} & Vincolo
	
	Obbligatorio & La parte client dovrà essere eseguibile in un browser HTML5 & Capitolato\tabularnewline
	\hline
	\begin{tikzpicture}
	\draw [->, thick] (0.2,0.2) -- (0.2,0.1) -- (1,0.1);
	\end{tikzpicture} & \hypertarget{R-3V3.1}{R-3V3.1} & Vincolo
	
	Obbligatorio & L'applicazione funziona corretamente su Chrome versione 47.0.2526 e successive & Capitolato\tabularnewline
	\hline
	\begin{tikzpicture}
	\draw [->, thick] (0.2,0.2) -- (0.2,0.1) -- (1,0.1);
	\end{tikzpicture} & \hypertarget{R-3V3.2}{R-3V3.2} & Vincolo
	
	Obbligatorio & L'applicazione funziona correttamente su Firefox versione 38.5.2esr e successive & Capitolato\tabularnewline
	\hline
	\begin{tikzpicture}
	\draw [->, thick] (0.2,0.2) -- (0.2,0.1) -- (1,0.1);
	\end{tikzpicture} & \hypertarget{R-3V3.3}{R-3V3.3} & Vincolo
	
	Obbligatorio & L'applicazione funziona correttamente su Microsoft Edge versione 25.10586 e successive & Capitolato\tabularnewline
	\hline
	\begin{tikzpicture}
	\draw [->, thick] (0.2,0.2) -- (0.2,0.1) -- (1,0.1);
	\end{tikzpicture} & \hypertarget{R-3V3.4}{R-3V3.4} & Vincolo
	
	Obbligatorio & L'applicazione funziona correttamente su Safari versione 9.0.2 e successive & Capitolato\tabularnewline
	\hline
	& \hypertarget{R-3V4}{R-3V4} & Vincolo
	
	Obbligatorio & L'applicazione deve utilizzare fogli stile CSS per l’ aspetto estetico & Capitolato\tabularnewline
	\hline
	& \hypertarget{R-3V5}{R-3V5} & Vincolo
	
	Obbligatorio & L'applicazione deve utilizzare Javascript per la parte attiva del client & Capitolato\tabularnewline
	\hline
	& \hypertarget{R-3V6}{R-3V6} & Vincolo
	
	Obbligatorio & La parte destinata agli esaminandi dovrà funzionare con qualunque dispositivo: dal Personal Computer, ai Tablet ed agli SmartPhone & Capitolato\tabularnewline
	\hline
	\begin{tikzpicture}
	\draw [->, thick] (0.2,0.2) -- (0.2,0.1) -- (1,0.1);
	\end{tikzpicture} & \hypertarget{R-3V6.1}{R-3V6.1} & Vincolo
	
	Obbligatorio & La parte destinata agli esaminandi dovrà funzionare si smartphone e tablet iOS dalla versione iOS6 & Capitolato\tabularnewline
	\hline
	\begin{tikzpicture}
	\draw [->, thick] (0.2,0.2) -- (0.2,0.1) -- (1,0.1);
	\end{tikzpicture} & \hypertarget{R-3V6.2}{R-3V6.2} & Vincolo
	
	Obbligatorio & La parte destinata agli esaminandi dovrà funzionare su Android dalla versione 4.0 & Capitolato\tabularnewline
	\hline
	& \hypertarget{R-3F7}{R-3F7} & Funzionale
	
	Obbligatorio & Il sistema deve gestire questionari & Capitolato\tabularnewline
	\hline
	\begin{tikzpicture}
	\draw [->, thick] (0.2,0.2) -- (0.2,0.1) -- (1,0.1);
	\end{tikzpicture} & \hypertarget{R-3F7.1}{R-3F7.1} & Funzionale
	
	Obbligatorio & L'applicazione deve archiviare quiz in un server, suddivisi per argomento
	& Capitolato\tabularnewline
	\hline
	\begin{tikzpicture}
	\draw [->, thick] (0.2,0.2) -- (0.2,0.1) -- (1,0.1);
	\end{tikzpicture} & \hypertarget{R-3F7.2}{R-3F7.2} & Funzionale
	
	Obbligatorio & L'applicazione deve poter archiviare questionari contenenti domande & Capitolato\tabularnewline
	\hline
	\begin{tikzpicture}
	\draw [->, thick] (0.2,0.2) -- (0.2,0.1) -- (1,0.1);
	\end{tikzpicture} & \hypertarget{R-3F7.3}{R-3F7.3} & Funzionale
	
	Obbligatorio & L'applicazione deve tradurre da QML a HTML le domande archiviate & Capitolato\tabularnewline
	\hline
	\begin{tikzpicture}
	\draw [->, thick] (0.2,0.2) -- (0.2,0.1) -- (1,0.1);
	\end{tikzpicture} & \hypertarget{R-3F7.4}{R-3F7.4} & Funzionale
	
	Obbligatorio & L'applicazione deve proporre questionari preconfezionati & Capitolato\tabularnewline
	\hline
	\begin{tikzpicture}
	\draw [->, thick] (0.2,0.2) -- (0.2,0.1) -- (1,0.1);
	\end{tikzpicture} & \hypertarget{R-3F7.5}{R-3F7.5} & Funzionale
	
	Obbligatorio & L'applicazione deve valutare le risposte fornite dall’ utente & Capitolato\tabularnewline
	\hline
	\begin{tikzpicture}
	\draw [->, thick] (0.2,0.2) -- (0.2,0.1) -- (1,0.1);
	\end{tikzpicture} & \hypertarget{R-3F7.6}{R-3F7.6} & Funzionale
	
	Obbligatorio & Le domande dovranno essere raccolte attraverso uno specifico linguaggio, che chiameremo Quiz Markup Language (QML), studiato per descrivere quiz
	& Capitolato\tabularnewline
	\hline
	\begin{tikzpicture}
	\draw [->, thick] (0.4,0.2) -- (0.4,0.1) -- (1,0.1);
	\end{tikzpicture} & \hypertarget{R-3F7.6.1}{R-3F7.6.1} & Funzionale
	
	Obbligatorio & Il QML deve gestire risposte vero/falso, risposte a scelta multipla, testi e immagini & Capitolato\tabularnewline
	\hline
	\begin{tikzpicture}
	\draw [->, thick] (0.4,0.2) -- (0.4,0.1) -- (1,0.1);
	\end{tikzpicture} & \hypertarget{R-3F7.6.2}{R-3F7.6.2} & Funzionale
	
	Obbligatorio & Le domande scritte in QML devono essere editabili dall'utente & Proponente\tabularnewline
	\hline
	\begin{tikzpicture}
	\draw [->, thick] (0.4,0.2) -- (0.4,0.1) -- (1,0.1);
	\end{tikzpicture} & \hypertarget{R-2F7.6.3}{R-2F7.6.3} & Funzionale
	
	Opzionale & Le domande possono essere insrite e editabili attraverso interfaccia grafica & Proponente\tabularnewline
	\hline
	\begin{tikzpicture}
	\draw [->, thick] (0.2,0.2) -- (0.2,0.1) -- (1,0.1);
	\end{tikzpicture} & \hypertarget{R-3F7.7}{R-3F7.7} & Funzionale
	
	Obbligatorio & Il sistema deve gestire sia questionari pubblici (esercitazione) che privati (esami) & Proponente\tabularnewline
	\hline
	\begin{tikzpicture}
	\draw [->, thick] (0.2,0.2) -- (0.2,0.1) -- (1,0.1);
	\end{tikzpicture} & \hypertarget{R-3F7.8}{R-3F7.8} & Funzionale
	
	Obbligatorio & I docenti costruiranno questionari scegliendo tra le domande archiviate nel server & Capitolato\tabularnewline
	\hline
	\begin{tikzpicture}
	\draw [->, thick] (0.2,0.2) -- (0.2,0.1) -- (1,0.1);
	\end{tikzpicture} & \hypertarget{R-2F7.9}{R-2F7.9} & Funzionale
	
	Opzionale & il sistema deve offrire la possibilità di creare dinamicamente i questionari & Capitolato\tabularnewline
	\hline
	\begin{tikzpicture}
	\draw [->, thick] (0.2,0.2) -- (0.2,0.1) -- (1,0.1);
	\end{tikzpicture} & \hypertarget{R-2F7.10}{R-2F7.10} & Funzionale
	
	Opzionale & Il l'applicazione deve essere predisposta ad un sistema di feedback per le domande & Capitolato\tabularnewline
	\hline
	\begin{tikzpicture}
	\draw [->, thick] (0.2,0.2) -- (0.2,0.1) -- (1,0.1);
	\end{tikzpicture} & \hypertarget{R-3F7.11}{R-3F7.11} & Funzionale
	
	Obbligatorio & Il sistema deve permettere la raccolta dei dati relativi a questionari svolti & Capitolato\tabularnewline
	\hline
	\begin{tikzpicture}
	\draw [->, thick] (0.4,0.2) -- (0.4,0.1) -- (1,0.1);
	\end{tikzpicture} & \hypertarget{R-3F7.11.1}{R-3F7.11.1} & Funzionale
	
	Obbligatorio & Il sistema deve permettere la rivisitazione del test svolto e dei risultati ottenuti & Capitolato\tabularnewline
	\hline
	\begin{tikzpicture}
	\draw [->, thick] (0.4,0.2) -- (0.4,0.1) -- (1,0.1);
	\end{tikzpicture} & \hypertarget{R-3F7.11.2}{R-3F7.11.2} & Funzionale
	
	Obbligatorio & Il sistema deve offrire la possibilità di analizzare i dati raccolti dai test effettuati & Capitolato\tabularnewline
	\hline
	\begin{tikzpicture}
	\draw [->, thick] (0.4,0.2) -- (0.4,0.1) -- (1,0.1);
	\end{tikzpicture} & \hypertarget{R-3F7.11.3}{R-3F7.11.3} & Funzionale
	
	Obbligatorio & Il sistema deve offrire la possibilità di visualizzazione sotto forma di grafico delle statistiche i dati raccolti & Capitolato\tabularnewline
	\hline
	& \hypertarget{R-3F8}{R-3F8} & Funzionale
	
	Obbligatorio & Il sistema deve permettere l’autenticazione
	di un utente & \hyperlink{UC2}{UC2}\tabularnewline
	\hline
	\begin{tikzpicture}
	\draw [->, thick] (0.2,0.2) -- (0.2,0.1) -- (1,0.1);
	\end{tikzpicture} & \hypertarget{R-3F8.1}{R-3F8.1} & Funzionale
	
	Obbligatorio & L’autenticazione richiede l’inserimento dell’e-mail e della password & \hyperlink{UC2}{UC2}\tabularnewline
	\hline
	\begin{tikzpicture}
	\draw [->, thick] (0.2,0.2) -- (0.2,0.1) -- (1,0.1);
	\end{tikzpicture} & \hypertarget{R-3F8.2}{R-3F8.2} & Funzionale
	
	Obbligatorio & Il sistema deve permettere la visualizzazione di un messaggio di errore in caso di email o password non valida & \hyperlink{UC2}{UC2}\tabularnewline
	\hline
	\begin{tikzpicture}
	\draw [->, thick] (0.2,0.2) -- (0.2,0.1) -- (1,0.1);
	\end{tikzpicture} & \hypertarget{R-2F8.3}{R-2F8.3} & Funzionale
	
	Opzionale & Il sistema deve permettere il recupero password & \hyperlink{UC2}{UC2}\tabularnewline
	\hline
	& \hypertarget{R-3F9}{R-3F9} & Funzionale
	
	Obbligatorio & Il candidato attraverso una specifica interfaccia potrà rispondere alle domande & Capitolato\tabularnewline
	\hline
	& \hypertarget{R-3F10}{R-3F10} & Funzionale
	
	Obbligatorio & Il sistema deve avere almeno 3 tipi di utenti: amministratore, docente e studente & Capitolato\tabularnewline
	\hline
	& \hypertarget{R-3F11}{R-3F11} & Funzionale
	
	Obbligatorio & Il sistema deve gestire la registrazione di un utente & \hyperlink{UC1}{UC1}\tabularnewline
	\hline
	\begin{tikzpicture}
	\draw [->, thick] (0.2,0.2) -- (0.2,0.1) -- (1,0.1);
	\end{tikzpicture} & \hypertarget{R-3F11.1}{R-3F11.1} & Funzionale
	
	Obbligatorio & Il sistema deve contenere già un account amministratore & Interno\tabularnewline
	\hline
	\begin{tikzpicture}
	\draw [->, thick] (0.2,0.2) -- (0.2,0.1) -- (1,0.1);
	\end{tikzpicture} & \hypertarget{R-3F11.2}{R-3F11.2} & Funzionale
	
	Obbligatorio & Il sistema deve permettere la registrazione
	di uno studente & \hyperlink{UC1}{UC1}\tabularnewline
	\hline
	\begin{tikzpicture}
	\draw [->, thick] (0.4,0.2) -- (0.4,0.1) -- (1,0.1);
	\end{tikzpicture} & \hypertarget{R-3F11.2.1}{R-3F11.2.1} & Funzionale
	
	Obbligatorio & La registrazione di uno studente richiede l’inserimento
	della mail, password, e facoltativamente nome e cognome & \hyperlink{UC1}{UC1}\tabularnewline
	\hline
	\begin{tikzpicture}
	\draw [->, thick] (0.4,0.2) -- (0.4,0.1) -- (1,0.1);
	\end{tikzpicture} & \hypertarget{R-3F11.2.2}{R-3F11.2.2} & Funzionale
	
	Obbligatorio & Il sistema deve permettere la visualizzazione di un messaggio di errore se la mail dello studente è già presente nel sistema & \hyperlink{UC1}{UC1}\tabularnewline
	\hline
	\begin{tikzpicture}
	\draw [->, thick] (0.4,0.2) -- (0.4,0.1) -- (1,0.1);
	\end{tikzpicture} & \hypertarget{R-2F11.2.3}{R-2F11.2.3} & Funzionale
	
	Opzionale & Il sistema deve permettere la visualizzazione di un messaggio di errore per password troppo corta & \hyperlink{UC1}{UC1}\tabularnewline
	\hline
	\begin{tikzpicture}
	\draw [->, thick] (0.2,0.2) -- (0.2,0.1) -- (1,0.1);
	\end{tikzpicture} & \hypertarget{R-3F11.3}{R-3F11.3} & Funzionale
	
	Obbligatorio & L'amministratore deve poter inviare un invito di registrazione ad un docente di cui conosce l'indirizzo email & \hyperlink{UC3.1}{UC3.1}\tabularnewline
	\hline
	\begin{tikzpicture}
	\draw [->, thick] (0.2,0.2) -- (0.2,0.1) -- (1,0.1);
	\end{tikzpicture} & \hypertarget{R-3F11.4}{R-3F11.4} & Funzionale
	
	Obbligatorio & Il sistema deve poter registrate un docente che ha ricevuto un invito di registrazione & \hyperlink{UC1}{UC1}\tabularnewline
	\hline
	\begin{tikzpicture}
	\draw [->, thick] (0.4,0.2) -- (0.4,0.1) -- (1,0.1);
	\end{tikzpicture} & \hypertarget{R-3F11.4.1}{R-3F11.4.1} & Funzionale
	
	Obbligatorio & La registrazione di un docente richiede l’inserimento della mail, password, e facoltativamente nome e cognome & \hyperlink{UC1}{UC1}\tabularnewline
	\hline
	\begin{tikzpicture}
	\draw [->, thick] (0.4,0.2) -- (0.4,0.1) -- (1,0.1);
	\end{tikzpicture} & \hypertarget{R-3F11.4.2}{R-3F11.4.2} & Funzionale
	
	Obbligatorio & Il sistema deve permettere la visualizzazione di un messaggio di errore se la mail del docente è già presente nel sistema & \hyperlink{UC1}{UC1}\tabularnewline
	\hline
	\begin{tikzpicture}
	\draw [->, thick] (0.4,0.2) -- (0.4,0.1) -- (1,0.1);
	\end{tikzpicture} & \hypertarget{R-3F11.4.3}{R-3F11.4.3} & Funzionale
	
	Obbligatorio & Il sistema deve permettere la visualizzazione di un messaggio di errore per password troppo corta & \hyperlink{UC1}{UC1}\tabularnewline
	\hline
	\begin{tikzpicture}
	\draw [->, thick] (0.2,0.2) -- (0.2,0.1) -- (1,0.1);
	\end{tikzpicture} & \hypertarget{R-3F11.5}{R-3F11.5} & Funzionale
	
	Obbligatorio & Il sistema deve permettere la creazione di un nuovo amministratore da parte di un altro amministratore & \hyperlink{UC3.4}{UC3.4}\tabularnewline
	\hline
	& \hypertarget{R-3V12}{R-3V12} & Vincolo
	
	Obbligatorio & Il database server side deve essere SQL (PostgreSQL) oppure noSQL (MongoDB) & Capitolato\tabularnewline
	\hline
	\caption{Tabella requisiti / fonti} \tabularnewline
	\end{longtable}


\newpage
\section{Use Cases}
\hypertarget{UC1}{}
\subsection{Caso d'uso UC1: Registrazione}
\begin{figure}[H]
	\centering

	\caption{Caso d'uso UC1: Registrazione}\label{fig:UC1} 
\end{figure}\begin{itemize}
\item \textbf{Attori}: Utente;
\item \textbf{Scopo e descrizione}: Utente non registrato; 
\item \textbf{Precondizione}: L’utente non è ancora registrato presso il sistema;

\item \textbf{Flusso principale degli eventi}:
\begin{enumerate}
	\item l'utente clicca sul tasto registrazione;
	\item inserisce un mail e password;
	\item l'utente può facoltativamente inserire Nome e Cognome;
	\item clicca sul tasto registrati;
	
\end{enumerate}
\item \textbf{Estensioni}:
\begin{enumerate}
	\item nel caso in cui l’utente inserisca una mail già censita dal sistema l’utente non viene registrato presso il sistema e viene visualizzato un errore esplicativo;
	\item nel caso in cui l’utente inserisca una password troppo corta l’utente non viene registrato presso il sistema e viene visualizzato un errore esplicativo;
	
\end{enumerate}
\item \textbf{Postcondizione}: L’utente possiede un’account di tipo studente o docente presso il sistema, contraddistinto da una username e da una password.
\end{itemize}
\hypertarget{UC2}{}
\subsection{Caso d'uso UC2: Autenticazione}
\begin{figure}[H]
	\centering

	\caption{Caso d'uso UC2: Autenticazione}\label{fig:UC2} 
\end{figure}\begin{itemize}
\item \textbf{Attori}: Utente;
\item \textbf{Scopo e descrizione}: Utente registrato; 
\item \textbf{Precondizione}: L'utente è registrato presso il sitema;

\item \textbf{Flusso principale degli eventi}:
\begin{enumerate}
	\item l'utente inserisce mail e password corretti;
	\item l'utente clicca su accedi;
	
\end{enumerate}
\item \textbf{Estensioni}:
\begin{enumerate}
	\item nel caso in cui l'utente inserisca una coppia di mail e password non presente nel sistema viene visualizzato un errore esplicativo e viene fornita la possibilità di ritentare l'accesso;
	
\end{enumerate}
\item \textbf{Postcondizione}: L'utente è autenticato presso il sistema.
\end{itemize}
\hypertarget{UC3}{}
\subsection{Caso d'uso UC3: Azioni Amministratore}
\begin{figure}[H]
	\centering

	\caption{Caso d'uso UC3: Azioni Amministratore}\label{fig:UC3} 
\end{figure}\begin{itemize}
\item \textbf{Attori}: Utente;
\item \textbf{Scopo e descrizione}: Amministratore; 
\item \textbf{Precondizione}: L'amministratore è autenticato nel sistema;

\item \textbf{Flusso principale degli eventi}:
\begin{enumerate}
	\item Amministratore può  invitare un docente (\hyperlink{UC3.1}{UC3.1});
	\item Amministratore può rimuovere un docente (\hyperlink{UC3.2}{UC3.2});
	\item Amministratore può rimuovere uno studente (\hyperlink{UC3.3}{UC3.3});
	\item Amministratore può creare un nuovo utente amministratore;
	\item Amministratore può gestire il proprio profilo (\hyperlink{UC9}{UC9});
	\item L'amministratore può effettuare il logout (\hyperlink{UC4}{UC4});
	
\end{enumerate}
\item \textbf{Postcondizione}: Il sistema ha ottenuto le informazioni sulle operazioni che
l’amministratore desidera eseguire.
\end{itemize}
\hypertarget{UC3.1}{}
\subsection{Caso d'uso UC3.1: Invita docente}
\begin{figure}[H]
	\centering

	\caption{Caso d'uso UC3.1: Invita docente}\label{fig:UC3.1} 
\end{figure}\begin{itemize}
\item \textbf{Attori}: Utente;
\item \textbf{Scopo e descrizione}: Amministratore; 
\item \textbf{Precondizione}: Amministratore è a conoscenza dell’email del docente da invitare;

\item \textbf{Flusso principale degli eventi}:
\begin{enumerate}
	\item Amministratore inserisce l'email del docente da invitare nel campo opportuno ;
	\item Amministratore invia al sistema l'email del docente da invitare;
	
\end{enumerate}
\item \textbf{Estensioni}:
\begin{enumerate}
	\item In caso di email già presente nel sistema viene segnalato all'Amministratore;
	
\end{enumerate}
\item \textbf{Postcondizione}: Sistema invia token di registrazione del docente alla medesima email.
\end{itemize}
\hypertarget{UC3.2}{}
\subsection{Caso d'uso UC3.2: Rimuovi  docente}
\begin{figure}[H]
	\centering

	\caption{Caso d'uso UC3.2: Rimuovi  docente}\label{fig:UC3.2} 
\end{figure}\begin{itemize}
\item \textbf{Attori}: Utente;
\item \textbf{Scopo e descrizione}: Amministratore; 
\item \textbf{Precondizione}: Il docente da rimuovere deve essere presente nel sistema;

\item \textbf{Flusso principale degli eventi}:
\begin{enumerate}
	\item Amministratore ricerca il docente da rimuovere;
	\item Amministratore seleziona il docente da rimuovere;
	\item Amministratore conferma la rimozione del docente selezionato;
	
\end{enumerate}
\item \textbf{Postcondizione}: Il medesimo docente è stato rimosso dal sistema.
\end{itemize}
\hypertarget{UC3.3}{}
\subsection{Caso d'uso UC3.3: Rimuovi  studente}
\begin{figure}[H]
	\centering

	\caption{Caso d'uso UC3.3: Rimuovi  studente}\label{fig:UC3.3} 
\end{figure}\begin{itemize}
\item \textbf{Attori}: Utente;
\item \textbf{Scopo e descrizione}: Amministratore; 
\item \textbf{Precondizione}: Lo studente da rimuovere deve essere presente nel sistema;

\item \textbf{Flusso principale degli eventi}:
\begin{enumerate}
	\item Amministratore ricerca lo studente da rimuovere;
	\item Amministratore seleziona lo studente da rimuovere;
	\item Amministratore conferma la rimozione dello studente selezionato;
	
\end{enumerate}
\item \textbf{Postcondizione}: Il medesimo studente è stato rimosso dal sistema.
\end{itemize}
\hypertarget{UC3.4}{}
\subsection{Caso d'uso UC3.4: Aggiunta amministratore}
\begin{figure}[H]
	\centering

	\caption{Caso d'uso UC3.4: Aggiunta amministratore}\label{fig:UC3.4} 
\end{figure}\begin{itemize}
\item \textbf{Attori}: Utente;
\item \textbf{Scopo e descrizione}: Amministratore; 
\item \textbf{Precondizione}: L'amministratore è autenticato nel sistema;

\item \textbf{Flusso principale degli eventi}:
\begin{enumerate}
	\item L'amministratore inserisce mail e password del nuovo amministratore;
	\item L'amministratore seleziona la funzionalità di aggiunta amministratore;
	
\end{enumerate}
\item \textbf{Estensioni}:
\begin{enumerate}
	\item nel caso in cui l’amministratore inserisca una mail già censita dal sistema l’utente non viene registrato presso il sistema e viene visualizzato un errore esplicativo;
	
\end{enumerate}
\item \textbf{Postcondizione}: È stato inserito un nuovo amministratore.
\end{itemize}
\hypertarget{UC4}{}
\subsection{Caso d'uso UC4: Logout}\begin{itemize}
	\item \textbf{Attori}: Utente;
	\item \textbf{Scopo e descrizione}: Utente Registrato; 
	\item \textbf{Precondizione}: L'utente è autenticato nel sistema;
	\item \textbf{Postcondizione}: L'utente non è più autenticato nel sistema.
\end{itemize}
\hypertarget{UC5}{}
\subsection{Caso d'uso UC5: Gestione classi}
\begin{figure}[H]
	\centering

	\caption{Caso d'uso UC5: Gestione classi}\label{fig:UC5} 
\end{figure}\begin{itemize}
\item \textbf{Attori}: Utente;
\item \textbf{Scopo e descrizione}: Docente; 
\item \textbf{Precondizione}: Il docente è autenticato nel sistema;

\item \textbf{Flusso principale degli eventi}:
\begin{enumerate}
	\item Il docente può creare una nuova classe (\hyperlink{UC5.1}{UC5.1});
	\item Il docente può modificare una classe (\hyperlink{UC5.2}{UC5.2});
	\item Il docente può eliminare una classe (\hyperlink{UC5.3}{UC5.3});
	
\end{enumerate}
\item \textbf{Postcondizione}: Il sistema ha ottenuto le informazioni sulle operazioni che il docente desidera eseguire sulla classe.
\end{itemize}
\hypertarget{UC5.1}{}
\subsection{Caso d'uso UC5.1: Inserisci classe}
\begin{figure}[H]
	\centering

	\caption{Caso d'uso UC5.1: Inserisci classe}\label{fig:UC5.1} 
\end{figure}\begin{itemize}
\item \textbf{Attori}: Utente;
\item \textbf{Scopo e descrizione}: Docente; 
\item \textbf{Precondizione}: Il docente è autenticato nel sistema;

\item \textbf{Flusso principale degli eventi}:
\begin{enumerate}
	\item Viene selezionata la funzionalità inserisci nuova classe;
	
\end{enumerate}
\item \textbf{Postcondizione}: È stato creata una nuova classe.
\end{itemize}
\hypertarget{UC5.2}{}
\subsection{Caso d'uso UC5.2: Modifica classe}\begin{itemize}
	\item \textbf{Attori}: Utente;
	\item \textbf{Scopo e descrizione}: Docente; 
	\item \textbf{Precondizione}: Il docente è autenticato nel sistema, è presente nel sistema la classe da modificare;
	\item \textbf{Postcondizione}: È stata modificata la classe.
\end{itemize}
\hypertarget{UC5.3}{}
\subsection{Caso d'uso UC5.3: Elimina classe}
\begin{figure}[H]
	\centering

	\caption{Caso d'uso UC5.3: Elimina classe}\label{fig:UC5.3} 
\end{figure}\begin{itemize}
\item \textbf{Attori}: Utente;
\item \textbf{Scopo e descrizione}: Docente; 
\item \textbf{Precondizione}: Il docente è autenticato nel sistema, è presente nel sistema la classe da eliminare;

\item \textbf{Flusso principale degli eventi}:
\begin{enumerate}
	\item Viene sezionata una classe;
	\item Viene sezionata la funzionalità elimina;
	\item Viene confermata l'eliminazione;
	
\end{enumerate}
\item \textbf{Postcondizione}: È stata eliminata la classe.
\end{itemize}
\hypertarget{UC6}{}
\subsection{Caso d'uso UC6: Gestione questionari}
\begin{figure}[H]
	\centering

	\caption{Caso d'uso UC6: Gestione questionari}\label{fig:UC6} 
\end{figure}\begin{itemize}
\item \textbf{Attori}: Utente;
\item \textbf{Scopo e descrizione}: Docente; 
\item \textbf{Precondizione}: Il docente è autenticato nel sistema;

\item \textbf{Flusso principale degli eventi}:
\begin{enumerate}
	\item Il docente può creare un nuovo questionario (\hyperlink{UC6.1}{UC6.1});
	\item Il docente può modificare un questionario (\hyperlink{UC6.2}{UC6.2});
	\item Il docente può eliminare un questionario (\hyperlink{UC6.3}{UC6.3});
	
\end{enumerate}
\item \textbf{Postcondizione}: Il sistema ha ottenuto le informazioni sulle operazioni che
il docente desidera eseguire su un questionario.
\end{itemize}
\hypertarget{UC6.1}{}
\subsection{Caso d'uso UC6.1: Inserisci questionario}
\begin{figure}[H]
	\centering

	\caption{Caso d'uso UC6.1: Inserisci questionario}\label{fig:UC6.1} 
\end{figure}\begin{itemize}
\item \textbf{Attori}: Utente;
\item \textbf{Scopo e descrizione}: Docente; 
\item \textbf{Precondizione}: Il docente è autenticato nel sistema;

\item \textbf{Flusso principale degli eventi}:
\begin{enumerate}
	\item Viene selezionata la funzionalità inserisci nuovo questionario;
	\item Vengono ricercate e selezionate le domande da inserire;
	\item Vengono inserite le domande nel questionario;
	\item Viene inserito l'argomento o gli argomenti del questionario;
	\item Viene confermato il questionario;
	
\end{enumerate}
\item \textbf{Postcondizione}: È stato creato un nuovo questionario.
\end{itemize}
\hypertarget{UC6.2}{}
\subsection{Caso d'uso UC6.2: Modifica questionario}
\begin{figure}[H]
	\centering

	\caption{Caso d'uso UC6.2: Modifica questionario}\label{fig:UC6.2} 
\end{figure}\begin{itemize}
\item \textbf{Attori}: Utente;
\item \textbf{Scopo e descrizione}: Docente; 
\item \textbf{Precondizione}: Il docente è autenticato nel sistema, è presente nel sistema il questionario da modificare;

\item \textbf{Flusso principale degli eventi}:
\begin{enumerate}
	\item Viene selezionata la funzionalità modifica questionario;
	\item Può venire cercata una domande da inserire filtrando per creatore, argomento, tipo domanda, testo domanda;
	\item Possono venire inserite domande al questionario;
	\item Possono essere eliminate dal questionario;
	\item Viene confermato il questionario;
	
\end{enumerate}
\item \textbf{Postcondizione}: È stato modificato il questionario.
\end{itemize}
\hypertarget{UC6.3}{}
\subsection{Caso d'uso UC6.3: Elimina questionario}
\begin{figure}[H]
	\centering

	\caption{Caso d'uso UC6.3: Elimina questionario}\label{fig:UC6.3} 
\end{figure}\begin{itemize}
\item \textbf{Attori}: Utente;
\item \textbf{Scopo e descrizione}: Docente; 
\item \textbf{Precondizione}: Il docente è autenticato nel sistema, è presente nel sistema il questionario da eliminare;

\item \textbf{Flusso principale degli eventi}:
\begin{enumerate}
	\item Viene selezionato un questionario;
	\item Viene selezionata la funzionalità elimina;
	\item Viene confermata l'eliminazione;
	
\end{enumerate}
\item \textbf{Postcondizione}: È stato eliminato il questionario.
\end{itemize}
\hypertarget{UC7}{}
\subsection{Caso d'uso UC7: Gestione domande}
\begin{figure}[H]
	\centering

	\caption{Caso d'uso UC7: Gestione domande}\label{fig:UC7} 
\end{figure}\begin{itemize}
\item \textbf{Attori}: Utente;
\item \textbf{Scopo e descrizione}: Docente; 
\item \textbf{Precondizione}: Il docente è autenticato nel sistema;

\item \textbf{Flusso principale degli eventi}:
\begin{enumerate}
	\item Il docente può creare una nuova domanda (\hyperlink{UC7.1}{UC7.1});
	\item Il docente può modificare una domanda (\hyperlink{UC7.2}{UC7.2});
	\item Il docente può eliminare una domanda (\hyperlink{UC7.3}{UC7.3});
	
\end{enumerate}
\item \textbf{Postcondizione}: Il sistema ha ottenuto le informazioni sulle operazioni che il docente desidera eseguire sulle domande.
\end{itemize}
\hypertarget{UC7.1}{}
\subsection{Caso d'uso UC7.1: Inserisci domanda}
\begin{figure}[H]
	\centering

	\caption{Caso d'uso UC7.1: Inserisci domanda}\label{fig:UC7.1} 
\end{figure}\begin{itemize}
\item \textbf{Attori}: Utente;
\item \textbf{Scopo e descrizione}: Docente; 
\item \textbf{Precondizione}: Il docente è autenticato nel sistema;

\item \textbf{Flusso principale degli eventi}:
\begin{enumerate}
	\item Viene selezionata da funzionalità inserisci domanda;
	
\end{enumerate}
\item \textbf{Postcondizione}: È stata creata una nuova domanda.
\end{itemize}
\hypertarget{UC7.2}{}
\subsection{Caso d'uso UC7.2: Modifica domanda}\begin{itemize}
	\item \textbf{Attori}: Utente;
	\item \textbf{Scopo e descrizione}: Docente; 
	\item \textbf{Precondizione}: Il docente è autenticato nel sistema, è presente nel sistema la domanda da modificare;
	\item \textbf{Postcondizione}: La domanda è stata modificata.
\end{itemize}
\hypertarget{UC7.3}{}
\subsection{Caso d'uso UC7.3: Elimina domanda}
\begin{figure}[H]
	\centering

	\caption{Caso d'uso UC7.3: Elimina domanda}\label{fig:UC7.3} 
\end{figure}\begin{itemize}
\item \textbf{Attori}: Utente;
\item \textbf{Scopo e descrizione}: Docente; 
\item \textbf{Precondizione}: Il docente è autenticato nel sistema, è presente nel sistema la domanda da eliminare;

\item \textbf{Flusso principale degli eventi}:
\begin{enumerate}
	\item Viene selezionata una domanda;
	\item Viene selezionata la funzionalità elimina;
	\item Viene confermata l'eliminazione;
	
\end{enumerate}
\item \textbf{Postcondizione}: La domanda è stata eliminata.
\end{itemize}
\hypertarget{UC8}{}
\subsection{Caso d'uso UC8: Esecuzione questionario}\begin{itemize}
	\item \textbf{Attori}: Utente;
	\item \textbf{Scopo e descrizione}: Studente, Docente; 
	\item \textbf{Precondizione}: L'utente è autenticato nel sistema;
	\item \textbf{Postcondizione}: Il questionario è stato compilato e vengono visualizzati i risultati.
\end{itemize}
\hypertarget{UC8.1}{}
\subsection{Caso d'uso UC8.1: Ricerca questionario}\begin{itemize}
	\item \textbf{Attori}: Utente;
	\item \textbf{Scopo e descrizione}: Studente, Docente; 
	\item \textbf{Precondizione}: L'utente è autenticato nel sistema;
	\item \textbf{Postcondizione}: Il sistema riceve in input i parametri di ricerca inseriti e mostra la lista dei questionari che li soddisfano.
\end{itemize}
\hypertarget{UC8.2}{}
\subsection{Caso d'uso UC8.2: Compilazione questionario}\begin{itemize}
	\item \textbf{Attori}: Utente;
	\item \textbf{Scopo e descrizione}: Studente, Docente; 
	\item \textbf{Precondizione}: L'utente è autenticato nel sistema; esiste il questionario da compilare;
	\item \textbf{Postcondizione}: Il questionario è stato compilato e vengono mostrati i risultati.
\end{itemize}
\hypertarget{UC9}{}
\subsection{Caso d'uso UC9: Gestione account}
\begin{figure}[H]
	\centering

	\caption{Caso d'uso UC9: Gestione account}\label{fig:UC9} 
\end{figure}\begin{itemize}
\item \textbf{Attori}: Utente;
\item \textbf{Scopo e descrizione}: Utente Autenticato; 
\item \textbf{Precondizione}: L’utente ha effettuato  il login con successo e decide di gestire il proprio account;

\item \textbf{Flusso principale degli eventi}:
\begin{enumerate}
	\item L’utente vede le proprie informazioni personali (\hyperlink{UC9.1}{UC9.1});
	\item L’utente modifica le proprie informazioni personali (\hyperlink{UC9.2}{UC9.2});
	
\end{enumerate}
\item \textbf{Postcondizione}:  Il sistema fornisce l’interfaccia di interazione per la gestione dell’account.
\end{itemize}
\hypertarget{UC9.1}{}
\subsection{Caso d'uso UC9.1: Visualizzazione informazioni account}
\begin{figure}[H]
	\centering

	\caption{Caso d'uso UC9.1: Visualizzazione informazioni account}\label{fig:UC9.1} 
\end{figure}\begin{itemize}
\item \textbf{Attori}: Utente;
\item \textbf{Scopo e descrizione}: L’utente prende visione di tutte le sue informazioni personali inserite al momento della
registrazione; 
\item \textbf{Precondizione}: Il sistema è pronto a fornire le informazioni personali dell’utente che ha fatto la richiesta;

\item \textbf{Flusso principale degli eventi}:
\begin{enumerate}
	\item L’utente visualizza la propria e-mail (\hyperlink{UC9.1.1}{UC9.1.1});
	\item L’utente visualizza il proprio nome (\hyperlink{UC9.1.2}{UC9.1.2});
	\item L’utente visualizza il proprio cognome (\hyperlink{UC9.1.3}{UC9.1.3});
	
\end{enumerate}
\item \textbf{Postcondizione}: Il sistema fornisce all’utente la possibilità di vedere le informazioni personali.
\end{itemize}
\hypertarget{UC9.1.1}{}
\subsection{Caso d'uso UC9.1.1:  Visualizzazione e-mail}\begin{itemize}
	\item \textbf{Attori}: Utente;
	\item \textbf{Scopo e descrizione}: L’utente visualizza il proprio indirizzo e-mail inserito al momento della registrazione; 
	\item \textbf{Precondizione}: L’utente ha richiesto di vedere il proprio indirizzo e-mail;
	\item \textbf{Postcondizione}: Il sistema ha fornito all’utente il dato richiesto.
\end{itemize}
\hypertarget{UC9.1.2}{}
\subsection{Caso d'uso UC9.1.2: Visualizzazione nome}\begin{itemize}
	\item \textbf{Attori}: Utente;
	\item \textbf{Scopo e descrizione}: L’utente visualizza il proprio nome inserito al momento della registrazione; 
	\item \textbf{Precondizione}: L’utente ha richiesto di vedere il proprio nome;
	\item \textbf{Postcondizione}: Il sistema ha fornito all’utente il dato richiesto.
\end{itemize}
\hypertarget{UC9.1.3}{}
\subsection{Caso d'uso UC9.1.3: Visualizzazione cognome}\begin{itemize}
	\item \textbf{Attori}: Utente;
	\item \textbf{Scopo e descrizione}: L’utente visualizza il proprio cognome inserito al momento della registrazione; 
	\item \textbf{Precondizione}: L’utente ha richiesto di vedere il proprio cognome;
	\item \textbf{Postcondizione}: Il sistema ha fornito all’utente il dato richiesto.
\end{itemize}
\hypertarget{UC9.2}{}
\subsection{Caso d'uso UC9.2: Modifica informazioni account}
\begin{figure}[H]
	\centering

	\caption{Caso d'uso UC9.2: Modifica informazioni account}\label{fig:UC9.2} 
\end{figure}\begin{itemize}
\item \textbf{Attori}: Utente;
\item \textbf{Scopo e descrizione}: L’utente vuole modificare le proprie generalità; 
\item \textbf{Precondizione}:  L’utente ha già preso visione delle proprie generalità ed ha deciso di fare alcune modifiche;

\item \textbf{Flusso principale degli eventi}:
\begin{enumerate}
	\item L’utente modifica il proprio nome (\hyperlink{UC9.2.1}{UC9.2.1});
	\item L’utente modifica il proprio cognome  (\hyperlink{UC9.2.2}{UC9.2.2});
	\item L’utente modifica la propria password (\hyperlink{UC9.2.3}{UC9.2.3});
	
\end{enumerate}
\item \textbf{Postcondizione}: Il sistema, dopo l’interazione con l’utente, salva le eventuali modifiche svolte.
\end{itemize}
\hypertarget{UC9.2.1}{}
\subsection{Caso d'uso UC9.2.1: Modifica nome}\begin{itemize}
	\item \textbf{Attori}: Utente;
	\item \textbf{Scopo e descrizione}: L’utente che ha visualizzato le proprie generalità, decide di voler cambiare il nome che
	ha inserito al momento della registrazione; 
	\item \textbf{Precondizione}: L’utente ha richiesto di modificare il proprio nome;
	\item \textbf{Postcondizione}: il sistema salva nome modificato.
\end{itemize}
\hypertarget{UC9.2.2}{}
\subsection{Caso d'uso UC9.2.2: Modifica cognome}\begin{itemize}
	\item \textbf{Attori}: Utente;
	\item \textbf{Scopo e descrizione}: L’utente che ha visualizzato le proprie generalità, decide di voler cambiare il cognome
	che ha inserito al momento della registrazione; 
	\item \textbf{Precondizione}: L’utente ha richiesto di modificare il proprio cognome;
	\item \textbf{Postcondizione}: il sistema salva  il cognome modificato .
\end{itemize}
\hypertarget{UC9.2.3}{}
\subsection{Caso d'uso UC9.2.3: Modifica password}
\begin{figure}[H]
	\centering

	\caption{Caso d'uso UC9.2.3: Modifica password}\label{fig:UC9.2.3} 
\end{figure}\begin{itemize}
\item \textbf{Scopo e descrizione}: L’utente che ha visualizzato le proprie generalità, decide di cambiare la password
inserendone una nuova; 
\item \textbf{Precondizione}: L’utente ha richiesto di modificare la propria password;

\item \textbf{Flusso principale degli eventi}:
\begin{enumerate}
	\item L’utente inserisce l’attuale password in uso (\hyperlink{UC9.2.3.1}{UC9.2.3.1});
	\item L’utente inserisce la nuova password (\hyperlink{UC9.2.3.2}{UC9.2.3.2});
	\item L’utente inserisce la password nuova una seconda volta (\hyperlink{UC9.2.3.3}{UC9.2.3.3});
	\item L’utente decide di confermare la modifica alla password di accesso (\hyperlink{UC9.2.3.4}{UC9.2.3.4});
	
\end{enumerate}
\item \textbf{Estensioni}:
\begin{enumerate}
	\item Errore input password (\hyperlink{UC9.2.4}{UC9.2.4});
	
\end{enumerate}
\item \textbf{Postcondizione}: il sistema salva  la password modificata.
\end{itemize}
\hypertarget{UC9.2.3.1}{}
\subsection{Caso d'uso UC9.2.3.1:  Inserimento vecchia password}\begin{itemize}
	\item \textbf{Attori}: Utente;
	\item \textbf{Scopo e descrizione}: L’utente ha deciso di modificare la password di accesso ed entra nella sezione per il cambio password; 
	\item \textbf{Precondizione}: Il sistema ha visualizzato la sezione dove poter inserire l’attuale password che andrà cambiata;
	\item \textbf{Postcondizione}: Il sistema convalida la password appena inserita.
\end{itemize}
\hypertarget{UC9.2.3.2}{}
\subsection{Caso d'uso UC9.2.3.2:  Inserimento nuova password}\begin{itemize}
	\item \textbf{Attori}: Utente;
	\item \textbf{Scopo e descrizione}: L’utente ha deciso di modificare la password di accesso ed entra nella sezione per il cambio password; 
	\item \textbf{Precondizione}: Il sistema ha visualizzato la sezione dove poter inserire la nuova password;
	\item \textbf{Postcondizione}: Il sistema convalida la nuova password appena inserita.
\end{itemize}
\hypertarget{UC9.2.3.3}{}
\subsection{Caso d'uso UC9.2.3.3: Conferma nuova password}\begin{itemize}
	\item \textbf{Attori}: Utente;
	\item \textbf{Scopo e descrizione}: L’utente ha deciso di modificare la password di accesso ed entra nella sezione per il cambio password; 
	\item \textbf{Precondizione}: Il sistema ha visualizzato la sezione dove inserire la conferma della nuova password;
	\item \textbf{Postcondizione}: Il sistema convalida la conferma della nuova password appena inserita.
\end{itemize}
\hypertarget{UC9.2.3.4}{}
\subsection{Caso d'uso UC9.2.3.4: Conferma modifica}\begin{itemize}
	\item \textbf{Attori}: Utente;
	\item \textbf{Scopo e descrizione}: L’utente autenticato deve poter confermare la modifica della password del proprio account; 
	\item \textbf{Precondizione}: Il sistema fornisce una schermata in cui è possibile confermare la modifica della password del proprio account;
	\item \textbf{Postcondizione}: L’utente ha confermato di voler modificare la propria password all’interno del sistema.
\end{itemize}
\hypertarget{UC9.2.4}{}
\subsection{Caso d'uso UC9.2.4: Visualizzazione messaggio errore password errata}
\begin{figure}[H]
	\centering

	\caption{Caso d'uso UC9.2.4: Visualizzazione messaggio errore password errata}\label{fig:UC9.2.4} 
\end{figure}\begin{itemize}
\item \textbf{Attori}: Utente;
\item \textbf{Scopo e descrizione}: L’utente che ha visualizzato le proprie informazioni personali, decide di voler cambiare la propria password, però non effettua la procedura di cambio password in modo corretto, quindi viene visualizzato il messaggio di errore; 
\item \textbf{Precondizione}: L’utente ha richiesto di modificare la password attuale;

\item \textbf{Flusso principale degli eventi}:
\begin{enumerate}
	\item L’utente vede il messaggio di errore per il cambiamento errato della nuova password;
	
\end{enumerate}
\item \textbf{Postcondizione}: Dopo la modifica della password, il sistema non salva il nuovo valore perché errato e visualizza un messaggio di errore.
\end{itemize}


\end{document}
